\subsection{string\_hash}\label{string_hash}
The \verb|string_hash| module is a hash map that maps short strings to an
\verb|int8_t| value. A key is as long as the name field \verb|line| of a single
time measurement (\verb|time_meas_t|, Section \ref{time}). A
typical use is looking up a handle by its name.

The library does not use \verb|malloc|: all nodes live in a fixed-size pool
inside the map structure. Where the map itself lives (global, as
\verb|static| or on the stack) is up to the caller. Memory footprint and run
time are therefore deterministic, and several independent maps are possible.

\subsubsection*{Configuration}
The sizes can be overridden before including \verb|string_hash.h| (or with
the compiler option \verb|-D|) (Listing \ref{sh_defines}).

\begin{listing}[H]
\begin{minted}{C}
#define SH_SIZE        100  // max. number of entries (node pool)
#define SH_BUCKET_CNT   37  // number of buckets, prime != 31
#define SH_KEY_LEN (TIME_MEAS_CHAR_PER_LINE + 1)  // incl. '\0'
\end{minted}
\caption{Configuration of \texttt{string\_hash}}
\label{sh_defines}
\end{listing}

A key may be at most \verb|SH_KEY_LEN - 1| characters long (20 characters
with \verb|TIME_MEAS_CHAR_PER_LINE = 20|); longer keys are rejected, not
truncated. \verb|SH_BUCKET_CNT| must not be a multiple of the hash multiplier
31, otherwise every character except the last drops out of the hash and the
entries are spread badly.

On the Cortex-M4 a node takes 28 bytes and a bucket 4 bytes. With the default
values a map therefore needs
$100 \cdot 28 + 37 \cdot 4 + 4 = 2952$\,bytes $\approx 2.9$\,kB of RAM.

\subsubsection*{State}
The map is described by \verb|sh_hash_map_t| (Listing \ref{sh_hash_map_t}).
\verb|bucket| holds the head of a linked list per bucket, \verb|node| is the
node pool. A node with \verb|used == false| is free. Collisions are resolved
by separate chaining; new nodes are inserted at the head of the list.

\begin{listing}[H]
\begin{minted}{C}
typedef struct sh_node_s {
    char key[SH_KEY_LEN];    // key, '\0'-terminated
    int8_t value;            // associated value
    bool used;               // node in use
    struct sh_node_s *next;  // next node in the bucket
} sh_node_t;

typedef struct sh_hash_map_s {
    int16_t numOfElements;            // number of entries
    sh_node_t *bucket[SH_BUCKET_CNT]; // list head per bucket
    sh_node_t node[SH_SIZE];          // node pool
} sh_hash_map_t;
\end{minted}
\caption{Definition of the \texttt{sh\_hash\_map\_t} structure}
\label{sh_hash_map_t}
\end{listing}

\subsubsection*{Hash function}
The bucket index is a polynomial hash over the characters of the key, computed
with Horner's scheme, modulo the number of buckets:
\[
h = \left(\sum_{i=0}^{n-1} k_i \cdot 31^{\,n-1-i}\right) \bmod
\mathtt{SH\_BUCKET\_CNT}
\]
The computation uses \verb|uint32_t|. Overflow on long keys is intended and
harmless, since the result only has to be evenly distributed, not exact.

\subsubsection*{API}
All functions check their pointer arguments and the key length. Invalid
arguments result in \verb|EM_ERR|.

{\sloppy\emergencystretch=3em
\begin{description}
\item[\texttt{sh\_init(mp)}] Resets the map: all buckets empty, all nodes
  free, \verb|numOfElements = 0|. Must be called before first use. Returns:
  \verb|em_msg| (\verb|EM_ERR| for a \verb|NULL| map, otherwise
  \verb|EM_OK|).
\item[\texttt{sh\_insert(mp, key, value)}] Inserts \verb|key| with
  \verb|value|. The key is copied into the node, so the caller need not keep
  it. If \verb|key| already exists, only its value is overwritten. Returns:
  \verb|em_msg| (\verb|EM_OK| after inserting or overwriting, \verb|EM_ERR|
  for an invalid argument, a key that is too long, or a full pool).
\item[\texttt{sh\_search(mp, key, value)}] Looks up \verb|key| and writes
  the associated value to \verb|*value|. Since every \verb|int8_t| is a valid
  value, ``not found'' is reported through the return value rather than
  a reserved value. Returns: \verb|em_msg| (\verb|EM_OK| if found;
  \verb|EM_ERR| if \verb|key| is missing or an argument is invalid,
  \verb|*value| is then left unchanged).
\item[\texttt{sh\_delete(mp, key)}] Removes \verb|key| from its list and
  returns the node to the pool. Returns: \verb|em_msg| (\verb|EM_OK| after
  removal, \verb|EM_ERR| if \verb|key| is missing or an argument is invalid).
\item[\texttt{sh\_hash\_function(key)}] Returns the bucket index of
  \verb|key|. Not needed for normal use, but helpful in tests. Returns:
  \verb|uint16_t| (index in the range
  $0 \ldots \mathtt{SH\_BUCKET\_CNT} - 1$).
\end{description}
\par}

\subsubsection*{Example}
\begin{listing}[H]
\begin{minted}{C}
static sh_hash_map_t names;
int8_t hdl;

sh_init(&names);
sh_insert(&names, "loop", 3);
if (sh_search(&names, "loop", &hdl) == EM_OK) {
    time_start(hdl, 0, NULL, NULL);
}
sh_delete(&names, "loop");
\end{minted}
\caption{Using \texttt{string\_hash}}
\label{sh_example}
\end{listing}

\subsubsection*{Limitations}
The map is not thread- or interrupt-safe; concurrent access from an ISR and
the main loop must be protected by the caller. Finding a free node is linear
in the pool ($O(\mathtt{SH\_SIZE})$); lookup and deletion are $O(1)$ on
average.
